% ch1-vector-spaces.tex — Chapter 1: Vector Spaces and Subspaces
% The foundational chapter of the course: vector space axioms, subspaces,
% linear combinations, bases, dimension, and the Grassmann formula.

\chapter{Vector Spaces and Subspaces}\label{ch:vector-spaces}

The concept of a \emph{vector space}\index{vector space} is the central
object of linear algebra. It abstracts the familiar geometry of arrows in
the plane and in space into a framework that encompasses polynomials,
matrices, functions, and solutions of differential equations --- all
governed by the same structural rules.

In the previous chapter we established the notion of a field~$\K$. A
vector space is, in essence, a set whose elements (called \emph{vectors})
can be added together and scaled by elements of~$\K$, subject to axioms
that generalize the intuitive properties of $\R^2$ and $\R^3$.

\medskip
\noindent\textbf{Why abstract?}\ \ The reader may wonder why we bother
with axioms when we could simply work in $\R^n$. The answer is twofold.
First, many naturally occurring mathematical objects---polynomial spaces,
function spaces, solution sets of linear systems---share the same
algebraic structure as $\R^n$, and it is both economical and illuminating
to develop the theory once, in the abstract, rather than repeat arguments
for each concrete example. Second, the abstract framework reveals
\emph{which properties matter}: the dimension of a space, for instance,
is a far more fundamental invariant than the specific nature of its
elements.

We begin with the motivating example of $\R^2$ and $\R^3$, then give the
full axiomatic definition and a rich collection of examples. The bulk of
the chapter develops the core structural theory: subspaces, linear
combinations, linear independence, bases, and dimension, culminating in
the Grassmann formula.

Throughout this chapter, $\K$ denotes a field (typically $\R$ or $\C$).


% ============================================================
\section{Geometric motivation}\label{sec:geometric-motivation}
% ============================================================

Before plunging into axioms, let us recall the geometric picture that
motivates the entire theory.

In the Euclidean plane $\R^2$, a vector $\vec{v} = (v_1, v_2)$ can be
visualized as an arrow from the origin to the point $(v_1, v_2)$. Two
fundamental operations are available:
\begin{itemize}
  \item \textbf{Addition:} $(v_1,v_2)+(w_1,w_2) = (v_1+w_1, v_2+w_2)$,
    which corresponds to the parallelogram rule.
  \item \textbf{Scalar multiplication:}
    $\lambda(v_1,v_2) = (\lambda v_1, \lambda v_2)$, which stretches or
    compresses the arrow (and reverses it if $\lambda < 0$).
\end{itemize}

These operations satisfy many natural properties: addition is commutative
and associative, there is a zero vector $(0,0)$, every vector has an
additive inverse, scalar multiplication distributes over addition, and so
on. The exact list of these properties constitutes the axioms of a vector
space.

\begin{figure}[ht]
\centering
\begin{tikzpicture}[>=Stealth, scale=1.2]
  % Grid
  \draw[gray!30, very thin, step=0.5] (-0.5,-0.5) grid (4.5,3.5);
  \draw[->] (-0.5,0) -- (4.5,0) node[right]{$x_1$};
  \draw[->] (0,-0.5) -- (0,3.5) node[above]{$x_2$};
  % Vectors
  \draw[->, thick, blue!70!black] (0,0) -- (2,1)
    node[midway, below right]{$\vec{v}$};
  \draw[->, thick, red!70!black] (0,0) -- (1,2)
    node[midway, above left]{$\vec{w}$};
  % Parallelogram
  \draw[dashed, blue!50] (2,1) -- (3,3);
  \draw[dashed, red!50] (1,2) -- (3,3);
  \draw[->, thick, purple!70!black] (0,0) -- (3,3)
    node[midway, above left]{$\vec{v}+\vec{w}$};
  % Scalar multiple
  \draw[->, thick, blue!40, dashed] (0,0) -- (3,1.5)
    node[right]{$\frac{3}{2}\vec{v}$};
\end{tikzpicture}
\caption{Vector addition (parallelogram rule) and scalar multiplication
in $\R^2$.}
\label{fig:vector-addition-R2}
\end{figure}

The same picture extends to $\R^3$, where vectors are triples
$(v_1,v_2,v_3)$ and the operations are defined componentwise. But what
about $\R^n$ for $n\geq 4$, where geometric visualization fails? And
what about objects like polynomials or matrices, which can also be added
and scaled? The axioms of a vector space capture exactly the common
structure.


% ============================================================
\section{Definition of a vector space}\label{sec:vs-definition}
% ============================================================

\begin{definition}{Vector space}{def:vector-space}
\index{vector space!definition}
A \emph{vector space over $\K$} (or \emph{$\K$-vector space}) is a set
$E$ equipped with two operations:
\begin{itemize}
  \item \textbf{Addition:}\index{addition!vector} a map
    $+\colon E\times E\to E$, $(u,v)\mapsto u+v$.
  \item \textbf{Scalar multiplication:}\index{scalar multiplication}
    a map $\cdot\colon\K\times E\to E$, $(\lambda,v)\mapsto \lambda\cdot v$
    (usually written $\lambda v$).
\end{itemize}
These operations satisfy the following axioms for all
$u,v,w\in E$ and all $\lambda,\mu\in\K$:

\medskip
\noindent\textbf{Addition axioms:}
\begin{enumerate}[label=\textbf{(V\arabic*)}]
  \item\label{ax:assoc} \textbf{Associativity:}
    $(u+v)+w = u+(v+w)$.
  \item\label{ax:comm} \textbf{Commutativity:}
    $u+v = v+u$.
  \item\label{ax:zero} \textbf{Zero vector:}\index{zero vector}
    there exists $\vzero\in E$ such that $v+\vzero = v$ for all $v\in E$.
  \item\label{ax:inverse} \textbf{Additive inverse:}
    for each $v\in E$ there exists $-v\in E$ such that $v+(-v)=\vzero$.
\end{enumerate}

\medskip
\noindent\textbf{Scalar multiplication axioms:}
\begin{enumerate}[label=\textbf{(V\arabic*)},resume]
  \item\label{ax:compat} \textbf{Compatibility:}
    $\lambda(\mu v) = (\lambda\mu) v$.
  \item\label{ax:unit} \textbf{Unit:}
    $1_\K\, v = v$.
  \item\label{ax:dist-scalar} \textbf{Distributivity over vector addition:}
    $\lambda(u+v) = \lambda u + \lambda v$.
  \item\label{ax:dist-field} \textbf{Distributivity over scalar addition:}
    $(\lambda+\mu)v = \lambda v + \mu v$.
\end{enumerate}
The elements of $E$ are called \emph{vectors}\index{vector} and the
elements of $\K$ are called \emph{scalars}\index{scalar}.
\end{definition}

\begin{remark}{Notation}{rem:notation-vs}
When the field $\K$ is clear from context we simply say ``vector space''
rather than ``$\K$-vector space''. We write $\vzero$ for the zero vector
and $0$ for the zero scalar, relying on context to distinguish them.
\end{remark}

\begin{proposition}{Elementary properties}{prop:elementary-vs}
\index{vector space!elementary properties}
Let $E$ be a $\K$-vector space. For all $v\in E$ and $\lambda\in\K$:
\begin{enumerate}[label=(\roman*)]
  \item The zero vector $\vzero$ is unique.
  \item The additive inverse $-v$ is unique.
  \item $0_\K\, v = \vzero$.
  \item $\lambda\,\vzero = \vzero$.
  \item $(-1)v = -v$.
  \item $\lambda v = \vzero$ implies $\lambda = 0$ or $v = \vzero$.
\end{enumerate}
\end{proposition}

\begin{proof}
\begin{enumerate}[label=(\roman*)]
  \item Suppose $\vzero$ and $\vzero'$ are both zero vectors. Then
    $\vzero = \vzero + \vzero' = \vzero'$, using the zero-vector
    property of $\vzero'$ and then of $\vzero$.

  \item Suppose $w$ and $w'$ are both additive inverses of $v$. Then
    $w = w + \vzero = w + (v + w') = (w + v) + w' = \vzero + w' = w'$.

  \item $0_\K\,v = (0_\K + 0_\K)v = 0_\K\,v + 0_\K\,v$ by~\ref{ax:dist-field}.
    Adding $-(0_\K\,v)$ to both sides gives $\vzero = 0_\K\,v$.

  \item $\lambda\,\vzero = \lambda(\vzero + \vzero) =
    \lambda\,\vzero + \lambda\,\vzero$ by~\ref{ax:dist-scalar}. Adding
    $-(\lambda\,\vzero)$ yields $\vzero = \lambda\,\vzero$.

  \item $v + (-1)v = 1\cdot v + (-1)v = (1+(-1))v = 0_\K\,v = \vzero$
    by~\ref{ax:dist-field} and (iii). By uniqueness of inverses,
    $(-1)v = -v$.

  \item If $\lambda\neq 0$, then $\lambda^{-1}$ exists in $\K$ and
    $v = 1\cdot v = (\lambda^{-1}\lambda)v = \lambda^{-1}(\lambda v) =
    \lambda^{-1}\vzero = \vzero$ by (iv). \qedhere
\end{enumerate}
\end{proof}


% ============================================================
\section{Key examples of vector spaces}\label{sec:vs-examples}
% ============================================================

We now present a gallery of examples that will accoSpany us throughout
the course. The reader is encouraged to verify the axioms in each case
(at least mentally).

% --- K^n ---
\begin{example}{The space $\K^n$}{ex:Kn}
\index{Kn@$\K^n$}
For any positive integer $n$, the set
\[
  \K^n = \setcond{(x_1,\ldots,x_n)}{x_1,\ldots,x_n\in\K}
\]
equipped with componentwise addition and scalar multiplication:
\begin{align*}
  (x_1,\ldots,x_n) + (y_1,\ldots,y_n) &= (x_1+y_1,\ldots,x_n+y_n), \\
  \lambda(x_1,\ldots,x_n) &= (\lambda x_1,\ldots,\lambda x_n),
\end{align*}
is a $\K$-vector space. The zero vector is $\vzero = (0,\ldots,0)$.

When $\K = \R$, the cases $n=1,2,3$ correspond to the real line, the
plane, and three-dimensional space, respectively.
\end{example}

% --- Matrices ---
\begin{example}{The space $\Mat_{n,p}(\K)$}{ex:matrix-space}
\index{matrix space}
The set $\Mat_{n,p}(\K)$ of all $n\times p$ matrices with entries in $\K$
is a $\K$-vector space under the usual matrix addition and scalar
multiplication. The zero vector is the $n\times p$ zero matrix.

By identifying an $n\times p$ matrix with an element of $\K^{np}$ (by
listing entries row by row), we see that $\Mat_{n,p}(\K)$ and $\K^{np}$
have the ``same'' structure; this will be made precise by the notion of
\emph{isomorphism} in 
ef{ch:linear-maps}.
\end{example}

% --- Polynomials ---
\begin{example}{Polynomial spaces}{ex:polynomial-space}
\index{polynomial space}
The set $\K[X]$ of all polynomials with coefficients in $\K$ is a
$\K$-vector space under the usual addition and scalar multiplication of
polynomials. The zero vector is the zero polynomial.

For each $n\in\N$, the subset
\[
  \K_n[X] \deff \setcond{P\in\K[X]}{\deg P\leq n}\cup\set{0}
\]
of polynomials of degree at most $n$ (together with the zero polynomial)
is also a $\K$-vector space.
\end{example}

% --- Function spaces ---
\begin{example}{Function spaces}{ex:function-space}
\index{function space}
Let $S$ be a nonempty set. The set $\K^S \deff \setcond{f}{f\colon S\to\K}$
of all functions from $S$ to $\K$ is a $\K$-vector space under pointwise
operations:
\begin{align*}
  (f+g)(s) &\deff f(s) + g(s), \\
  (\lambda f)(s) &\deff \lambda\cdot f(s),
\end{align*}
for all $s\in S$, $\lambda\in\K$. The zero vector is the constant
function $s\mapsto 0$.

Important special cases include:
\begin{itemize}
  \item $\cC^0([a,b],\R)$: the space of continuous real-valued functions
    on $[a,b]$.
  \item $\cC^\infty(\R,\R)$: the space of infinitely differentiable
    functions.
  \item The space of sequences $(\K^\N)$: functions from $\N$ to $\K$.
\end{itemize}
\end{example}

% --- Solution spaces ---
\begin{example}{Solution space of a homogeneous linear system}{ex:homogeneous-system}
\index{solution space}\index{homogeneous system}
Let $A\in\Mat_{n,p}(\K)$. The set of solutions
\[
  \cS = \setcond{X\in\K^p}{AX = \vzero}
\]
of the homogeneous system $AX = \vzero$ is a $\K$-vector space (a
subspace of $\K^p$, as we shall see in~
ef{sec:subspaces}).

Indeed, if $AX_1 = \vzero$ and $AX_2 = \vzero$, then
$A(X_1+X_2) = AX_1 + AX_2 = \vzero$ and
$A(\lambda X_1) = \lambda AX_1 = \lambda\vzero = \vzero$.

\textbf{Warning:} The solution set of a \emph{non}-homogeneous system
$AX = B$ (with $B\neq\vzero$) is \emph{not} a vector space, since it
does not contain $\vzero$ (unless $B=\vzero$).
\end{example}

\begin{example}{The trivial vector space}{ex:trivial-space}
\index{trivial vector space}
The set $\set{\vzero}$, consisting of a single element, is a
$\K$-vector space (the only possibility for the operations being
$\vzero+\vzero = \vzero$ and $\lambda\,\vzero = \vzero$). It is called
the \emph{trivial} or \emph{zero} vector space.
\end{example}

\begin{remark}{The field $\K$ is a vector space over itself}{rem:K-over-K}
The field $\K$ is itself a $\K$-vector space (with field multiplication
playing the role of scalar multiplication). More generally, $\C$ is an
$\R$-vector space (of dimension~$2$, as we shall see).
\end{remark}


% ============================================================
\section{Subspaces}\label{sec:subspaces}
% ============================================================

A subspace is a subset of a vector space that is itself a vector space
(with the inherited operations). Rather than checking all eight axioms,
the following characterization provides a convenient shortcut.

\begin{definition}{Vector subspace}{def:subspace}
\index{subspace!definition}
Let $E$ be a $\K$-vector space. A subset $F\subseteq E$ is a
\emph{vector subspace} (or simply \emph{subspace}) of $E$ if:
\begin{enumerate}[label=(\roman*)]
  \item $F$ is nonempty (equivalently, $\vzero\in F$).
  \item $F$ is closed under addition: $u+v\in F$ for all $u,v\in F$.
  \item $F$ is closed under scalar multiplication: $\lambda v\in F$ for
    all $\lambda\in\K$ and $v\in F$.
\end{enumerate}
\end{definition}

\begin{proposition}{Subspace characterization}{prop:subspace-char}
\index{subspace!characterization}
A subset $F\subseteq E$ is a subspace of $E$ if and only if
\begin{enumerate}[label=(\alph*)]
  \item $F\neq\varnothing$ (or equivalently, $\vzero\in F$), and
  \item for all $\lambda\in\K$ and all $u,v\in F$:
    $\lambda u + v\in F$.
\end{enumerate}
Equivalently, $F$ is nonempty and closed under linear combinations: for
all $\lambda,\mu\in\K$ and $u,v\in F$, $\lambda u + \mu v\in F$.
\end{proposition}

\begin{proof}
If $F$ is a subspace, then (a) and (b) are immediate from the definition.
Conversely, assume (a) and (b). Taking $\lambda = 0$ in (b) shows
$\vzero + v = v\in F$ for any $v\in F$, confirming closure under
addition (take $\lambda=1$). Taking $v = \vzero$ in (b) shows
$\lambda u\in F$, confirming closure under scalar multiplication. Thus
$F$ satisfies all three conditions of the definition.
\end{proof}

\begin{remark}{The nonemptiness condition}{rem:nonempty}
Condition~(a) cannot be dropped. The empty set $\varnothing$ is trivially
closed under addition and scalar multiplication (vacuously) but is not a
vector space.
\end{remark}

\begin{example}{Subspaces of $\R^2$}{ex:subspaces-R2}
\index{subspace!of $\R^2$}
The subspaces of $\R^2$ are precisely:
\begin{enumerate}[label=(\roman*)]
  \item $\set{\vzero}$ (the origin).
  \item Lines through the origin: $\setcond{t\vec{v}}{t\in\R}$ for
    some $\vec{v}\neq\vzero$.
  \item $\R^2$ itself.
\end{enumerate}
Note that a line \emph{not} passing through the origin (e.g.\
$\setcond{(x,y)\in\R^2}{x+y=1}$) is \emph{not} a subspace.
\end{example}

\begin{figure}[ht]
\centering
\begin{tikzpicture}[>=Stealth, scale=1.0]
  % Axes
  \draw[gray!30, very thin, step=1] (-3.2,-2.2) grid (3.2,2.7);
  \draw[->] (-3.2,0) -- (3.2,0) node[right]{$x_1$};
  \draw[->] (0,-2.2) -- (0,2.7) node[above]{$x_2$};
  % Subspace: line through origin
  \draw[thick, blue!70!black] (-3,-1.5) -- (3,1.5);
  \node[blue!70!black, above left] at (2.4,1.2)
    {$F = \Span(\vec{v})$};
  \draw[->, very thick, blue!70!black] (0,0) -- (2,1)
    node[below right]{$\vec{v}$};
  % NOT a subspace: translated line
  \draw[thick, red!70!black, dashed] (-2.5,2.5+0.75) -- (2.5,-2.5+0.75);
  \node[red!70!black, right] at (1.8,-0.75)
    {$x_1+x_2=1$ (not a subspace)};
  % Origin
  \fill (0,0) circle (2pt) node[below left]{$\vzero$};
\end{tikzpicture}
\caption{A subspace of $\R^2$ (a line through the origin) versus a
non-subspace (a translated line).}
\label{fig:subspace-R2}
\end{figure}

\begin{example}{Subspaces of $\R^3$}{ex:subspaces-R3}
\index{subspace!of $\R^3$}
The subspaces of $\R^3$ are:
\begin{enumerate}[label=(\roman*)]
  \item $\set{\vzero}$.
  \item Lines through the origin.
  \item Planes through the origin (e.g.\
    $\setcond{(x,y,z)\in\R^3}{ax+by+cz=0}$ with $(a,b,c)\neq\vzero$).
  \item $\R^3$ itself.
\end{enumerate}
This classification will follow from the theory of dimension developed
later in this chapter.
\end{example}

\begin{example}{Further examples of subspaces}{ex:further-subspaces}
\begin{enumerate}[label=(\alph*)]
  \item The set of symmetric $n\times n$ matrices
    $\mathrm{Sym}_n(\K) = \setcond{A\in\Mat_n(\K)}{\transpose{A}=A}$ is a
    subspace of $\Mat_n(\K)$.
  \item The set of skew-symmetric matrices
    $\mathrm{Alt}_n(\K) = \setcond{A\in\Mat_n(\K)}{\transpose{A}=-A}$
    is a subspace of $\Mat_n(\K)$.
  \item The set of polynomials of degree at most $n$, $\K_n[X]$, is a
    subspace of $\K[X]$.
  \item The set of even polynomials
    $\setcond{P\in\K[X]}{P(-X)=P(X)}$ is a subspace of $\K[X]$.
  \item The solution set of the differential equation $y'' + y = 0$ (a
    subset of $\cC^\infty(\R,\R)$) is a subspace (the zero function is a
    solution, and linear combinations of solutions are solutions).
\end{enumerate}
\end{example}


% ============================================================
\section{Intersection and sum of subspaces}\label{sec:intersection-sum}
% ============================================================

\begin{proposition}{Intersection of subspaces}{prop:intersection-subspace}
\index{intersection of subspaces}
Let $(F_i)_{i\in I}$ be a (possibly infinite) family of subspaces of a
$\K$-vector space $E$. Then $\displaystyle\bigcap_{i\in I} F_i$ is a
subspace of $E$.
\end{proposition}

\begin{proof}
Let $F = \bigcap_{i\in I} F_i$.
\begin{itemize}
  \item \textbf{Nonempty:} $\vzero\in F_i$ for every $i$, so
    $\vzero\in F$.
  \item \textbf{Closed under linear combinations:} Let $u,v\in F$ and
    $\lambda,\mu\in\K$. For each $i\in I$, $u,v\in F_i$ and $F_i$ is a
    subspace, so $\lambda u+\mu v\in F_i$. Hence
    $\lambda u+\mu v\in F$. \qedhere
\end{itemize}
\end{proof}

\begin{remark}{Union is not a subspace in general}{rem:union-not-subspace}
\index{union of subspaces}
In contrast, the \emph{union} of two subspaces is generally \emph{not}
a subspace. For instance, in $\R^2$, let $F_1$ be the $x$-axis and $F_2$
the $y$-axis. Then $(1,0)\in F_1$ and $(0,1)\in F_2$, but
$(1,0)+(0,1)=(1,1)\notin F_1\cup F_2$.

In fact, $F_1\cup F_2$ is a subspace if and only if $F_1\subseteq F_2$
or $F_2\subseteq F_1$.
\end{remark}

The correct replacement for union is the \emph{sum} of subspaces.

\begin{definition}{Sum of subspaces}{def:sum-subspaces}
\index{sum of subspaces}
Let $F_1,\ldots,F_r$ be subspaces of $E$. Their \emph{sum} is
\[
  F_1 + F_2 + \cdots + F_r \deff
  \setcond{v_1 + v_2 + \cdots + v_r}{v_i\in F_i \text{ for each } i}.
\]
For two subspaces, $F+G = \setcond{u+v}{u\in F,\;v\in G}$.
\end{definition}

\begin{proposition}{Sum is the smallest containing subspace}{prop:sum-smallest}
The sum $F_1+\cdots+F_r$ is a subspace of $E$. It is the smallest
subspace of $E$ containing each $F_i$, i.e.\
$F_1+\cdots+F_r = \bigcap\setcond{H}{H \text{ subspace of } E,\;
F_i\subseteq H \text{ for all } i}$.
\end{proposition}

\begin{proof}
We prove the case $r=2$; the general case follows by induction.

Let $S = F + G$. First, $\vzero = \vzero + \vzero\in S$, so $S$ is
nonempty. Let $u_1+v_1, u_2+v_2\in S$ (with $u_i\in F$, $v_i\in G$) and
$\lambda\in\K$. Then
\[
  \lambda(u_1+v_1)+(u_2+v_2) = (\lambda u_1+u_2)+(\lambda v_1+v_2)\in S,
\]
since $\lambda u_1+u_2\in F$ and $\lambda v_1+v_2\in G$. Hence $S$ is a
subspace.

For the minimality claim: $F\subseteq S$ (take $v=\vzero$) and
$G\subseteq S$ (take $u=\vzero$), so $S$ contains both. If $H$ is any
subspace containing $F$ and $G$, then for any $u+v\in S$ with $u\in F$
and $v\in G$, we have $u,v\in H$, so $u+v\in H$; thus $S\subseteq H$.
\end{proof}


% ============================================================
\section{Direct sum and complementary subspaces}\label{sec:direct-sum}
% ============================================================

\begin{definition}{Direct sum (internal)}{def:direct-sum}
\index{direct sum}\index{direct sum!internal}
Let $F$ and $G$ be subspaces of $E$. We say that $E$ is the
\emph{direct sum} of $F$ and $G$, written $E = F\oplus G$, if:
\begin{enumerate}[label=(\roman*)]
  \item $E = F + G$ (every vector in $E$ can be written as $u+v$ with
    $u\in F$, $v\in G$), and
  \item $F\cap G = \set{\vzero}$.
\end{enumerate}
Equivalently, every $x\in E$ can be written \emph{uniquely} as $x=u+v$
with $u\in F$ and $v\in G$.
\end{definition}

\begin{proposition}{Characterization of direct sum}{prop:direct-sum-char}
\index{direct sum!characterization}
Let $F$ and $G$ be subspaces of $E$ with $E = F+G$. The following are
equivalent:
\begin{enumerate}[label=(\roman*)]
  \item $E = F\oplus G$.
  \item $F\cap G = \set{\vzero}$.
  \item Every $x\in E$ has a \emph{unique} decomposition $x = u + v$
    with $u\in F$ and $v\in G$.
\end{enumerate}
\end{proposition}

\begin{proof}
\textbf{(i)$\Leftrightarrow$(ii)} is immediate from the definition.

\textbf{(ii)$\Rightarrow$(iii):} Suppose $x = u_1+v_1 = u_2+v_2$ with
$u_i\in F$ and $v_i\in G$. Then $u_1-u_2 = v_2-v_1\in F\cap G = \set{\vzero}$,
so $u_1=u_2$ and $v_1=v_2$.

\textbf{(iii)$\Rightarrow$(ii):} The zero vector can be written as
$\vzero = \vzero + \vzero$. If $w\in F\cap G$, then
$\vzero = w + (-w)$ with $w\in F$ and $-w\in G$. By uniqueness,
$w = \vzero$.
\end{proof}

\begin{definition}{Complementary subspace}{def:complement}
\index{complementary subspace}
If $E = F\oplus G$, we say that $G$ is a \emph{complement} of $F$ in
$E$ (and $F$ is a complement of $G$). In this case, $F$ and $G$ are
called \emph{complementary subspaces}.
\end{definition}

\begin{remark}{Non-uniqueness of complements}{rem:complement-nonunique}
A subspace may have many different complements. For instance, in $\R^2$,
the $x$-axis has any non-horizontal line through the origin as a
complement.
\end{remark}

The direct sum extends naturally to more than two subspaces.

\begin{definition}{Direct sum of $r$ subspaces}{def:direct-sum-r}
\index{direct sum!of multiple subspaces}
Let $F_1,\ldots,F_r$ be subspaces of $E$. We say that
$E = F_1\oplus\cdots\oplus F_r$ if:
\begin{enumerate}[label=(\roman*)]
  \item $E = F_1 + \cdots + F_r$, and
  \item for each $j\in\set{1,\ldots,r}$,
    $F_j \cap (F_1+\cdots+F_{j-1}+F_{j+1}+\cdots+F_r) = \set{\vzero}$.
\end{enumerate}
Equivalently, every $x\in E$ has a unique decomposition
$x = v_1 + \cdots + v_r$ with $v_i\in F_i$.
\end{definition}

\begin{example}{Direct sum decomposition in $\R^3$}{ex:direct-sum-R3}
\index{direct sum!example in $\R^3$}
Let
\begin{align*}
  F &= \setcond{(x,0,0)}{x\in\R} &&\text{(the $x$-axis)}, \\
  G &= \setcond{(0,y,z)}{y,z\in\R} &&\text{(the $yz$-plane)}.
\end{align*}
Then $\R^3 = F\oplus G$, since every $(x,y,z) = (x,0,0)+(0,y,z)$
uniquely and $F\cap G = \set{\vzero}$.
\end{example}

\begin{figure}[ht]
\centering
\begin{tikzpicture}[>=Stealth, scale=1.1,
  x={(-0.4cm,-0.3cm)}, y={(1cm,0cm)}, z={(0cm,1cm)}]
  % Plane G (yz-plane)
  \fill[blue!10, opacity=0.6] (0,-2.2,-1.8) -- (0,2.2,-1.8) --
    (0,2.2,1.8) -- (0,-2.2,1.8) -- cycle;
  \node[blue!70!black] at (0,1.8,1.5) {$G$};
  % Line F (x-axis)
  \draw[thick, red!70!black] (2.8,0,0) -- (-2.8,0,0);
  \node[red!70!black, above] at (-2.2,0,0) {$F$};
  % Axes
  \draw[->] (0,0,0) -- (-3,0,0) node[left]{$x$};
  \draw[->] (0,0,0) -- (0,2.5,0) node[right]{$y$};
  \draw[->] (0,0,0) -- (0,0,2.2) node[above]{$z$};
  % Point and decomposition
  \coordinate (P) at (-1.5,1.2,0.8);
  \coordinate (Pf) at (-1.5,0,0);
  \coordinate (Pg) at (0,1.2,0.8);
  \fill[black] (P) circle (1.5pt) node[above right]{$(x,y,z)$};
  \draw[dashed, red!50] (P) -- (Pg);
  \draw[dashed, blue!50] (P) -- (Pf);
  \draw[->, thick, red!70!black] (0,0,0) -- (Pf)
    node[below]{$(x,0,0)$};
  \draw[->, thick, blue!70!black] (0,0,0) -- (Pg)
    node[right]{$(0,y,z)$};
  \fill (0,0,0) circle (1.5pt);
\end{tikzpicture}
\caption{Direct sum decomposition $\R^3 = F\oplus G$ where $F$ is the
$x$-axis and $G$ is the $yz$-plane.}
\label{fig:direct-sum-R3}
\end{figure}


% ============================================================
\section{Linear combinations and Span}\label{sec:Span}
% ============================================================

\begin{definition}{Linear combination}{def:linear-combination}
\index{linear combination}
Let $E$ be a $\K$-vector space and $v_1,\ldots,v_n\in E$. A
\emph{linear combination} of $v_1,\ldots,v_n$ is any vector of the form
\[
  \lambda_1 v_1 + \lambda_2 v_2 + \cdots + \lambda_n v_n,
  \qquad \lambda_1,\ldots,\lambda_n\in\K.
\]
The scalars $\lambda_1,\ldots,\lambda_n$ are called the
\emph{coefficients} of the linear combination.
\end{definition}

\begin{definition}{Span (generating set)}{def:Span}
\index{Span}\index{generating family}
Let $\cF = (v_1,\ldots,v_n)$ be a family of vectors in $E$. The
\emph{Span} of $\cF$ is
\[
  \Span(v_1,\ldots,v_n) \deff
  \setcond{\lambda_1 v_1 + \cdots + \lambda_n v_n}
    {\lambda_1,\ldots,\lambda_n\in\K}.
\]
It is the set of all linear combinations of $v_1,\ldots,v_n$.

A family $\cF$ is a \emph{generating family} (or \emph{Spanning set})
for $E$ if $\Span(\cF) = E$. In this case we say that $\cF$
\emph{Spans} (or \emph{generates}) $E$.
\end{definition}

\begin{proposition}{Span is a subspace}{prop:Span-subspace}
$\Span(v_1,\ldots,v_n)$ is a subspace of $E$. Moreover, it is the
smallest subspace containing $v_1,\ldots,v_n$.
\end{proposition}

\begin{proof}
Let $W = \Span(v_1,\ldots,v_n)$.
\begin{itemize}
  \item $\vzero = 0v_1+\cdots+0v_n\in W$, so $W\neq\varnothing$.
  \item If $u = \sum_i \alpha_i v_i$ and $w = \sum_i\beta_i v_i$ are in
    $W$ and $\lambda\in\K$, then
    $\lambda u + w = \sum_i(\lambda\alpha_i+\beta_i)v_i\in W$.
\end{itemize}
Hence $W$ is a subspace. It contains each $v_j$ (take
$\lambda_j=1$ and $\lambda_i=0$ for $i\neq j$). If $H$ is any subspace
containing every $v_j$, then $H$ contains all linear combinations of
$v_1,\ldots,v_n$, so $W\subseteq H$.
\end{proof}

\begin{example}{Spanning $\R^3$}{ex:Span-R3}
The canonical vectors $\vece{1}=(1,0,0)$, $\vece{2}=(0,1,0)$,
$\vece{3}=(0,0,1)$ Span $\R^3$, since every $(a,b,c)\in\R^3$ equals
$a\,\vece{1}+b\,\vece{2}+c\,\vece{3}$.
\end{example}

\begin{example}{Spanning a plane in $\R^3$}{ex:Span-plane}
The vectors $\vec{u}=(1,1,0)$ and $\vec{v}=(0,1,1)$ Span the plane
\[
  \Span(\vec{u},\vec{v}) = \setcond{(s,s+t,t)}{s,t\in\R}
  = \setcond{(x,y,z)\in\R^3}{x-y+z=0}.
\]
\end{example}


% ============================================================
\section{Linear independence}\label{sec:linear-independence}
% ============================================================

\begin{definition}{Linear independence}{def:linear-independence}
\index{linear independence}
A family $(v_1,\ldots,v_n)$ of vectors in $E$ is \emph{linearly
independent} (or \emph{free}) if
\[
  \lambda_1 v_1 + \cdots + \lambda_n v_n = \vzero
  \implies
  \lambda_1 = \lambda_2 = \cdots = \lambda_n = 0.
\]
A family that is not linearly independent is called \emph{linearly
dependent}\index{linear dependence}.
\end{definition}

\begin{remark}{Interpretation}{rem:li-interpretation}
Linear independence means that no vector in the family can be written as
a linear combination of the others. Equivalently, the only way to
represent $\vzero$ as a linear combination of $v_1,\ldots,v_n$ is the
\emph{trivial} one (all coefficients zero).
\end{remark}

\begin{proposition}{Characterization of linear dependence}{prop:ld-char}
\index{linear dependence!characterization}
A family $(v_1,\ldots,v_n)$ with $n\geq 2$ is linearly dependent if and
only if at least one vector $v_j$ can be written as a linear combination
of the others:
\[
  v_j = \sum_{\substack{i=1\\i\neq j}}^n \mu_i\, v_i.
\]
\end{proposition}

\begin{proof}
$(\Rightarrow)$~If the family is dependent, there exist
$\lambda_1,\ldots,\lambda_n$, not all zero, with
$\sum_i\lambda_i v_i = \vzero$. Pick $j$ with $\lambda_j\neq 0$. Then
$v_j = -\lambda_j^{-1}\sum_{i\neq j}\lambda_i v_i$.

$(\Leftarrow)$~If $v_j = \sum_{i\neq j}\mu_i v_i$, then
$\sum_{i\neq j}\mu_i v_i + (-1)v_j = \vzero$ is a nontrivial relation.
\end{proof}

\begin{proposition}{Properties of linear independence}{prop:li-properties}
\index{linear independence!properties}
Let $E$ be a $\K$-vector space.
\begin{enumerate}[label=(\roman*)]
  \item A single nonzero vector $(v)$ is linearly independent.
  \item Any subfamily of a linearly independent family is linearly
    independent.
  \item If $(v_1,\ldots,v_n)$ is linearly independent and $w\notin
    \Span(v_1,\ldots,v_n)$, then $(v_1,\ldots,v_n,w)$ is linearly
    independent.
  \item Any family containing $\vzero$ is linearly dependent.
  \item Any family containing a repeated vector is linearly dependent.
\end{enumerate}
\end{proposition}

\begin{proof}
\begin{enumerate}[label=(\roman*)]
  \item If $\lambda v = \vzero$ with $v\neq\vzero$, then $\lambda = 0$
    by 
ef{prop:prop:elementary-vs}(vi).

  \item Suppose $(v_{i_1},\ldots,v_{i_k})$ is a subfamily and
    $\sum_{j=1}^k \lambda_j v_{i_j} = \vzero$. Extend to a combination
    of the full family by setting the remaining coefficients to zero;
    independence of the full family forces all $\lambda_j = 0$.

  \item Suppose $\lambda_1 v_1+\cdots+\lambda_n v_n + \mu w = \vzero$.
    If $\mu\neq 0$, then
    $w = -\mu^{-1}(\lambda_1 v_1+\cdots+\lambda_n v_n)\in
    \Span(v_1,\ldots,v_n)$, contradicting the hypothesis. So $\mu=0$,
    and then $\lambda_1=\cdots=\lambda_n=0$ by independence.

  \item $1\cdot\vzero = \vzero$ is a nontrivial relation.

  \item If $v_i = v_j$ for $i\neq j$, then
    $1\cdot v_i + (-1)\cdot v_j = \vzero$ is nontrivial. \qedhere
\end{enumerate}
\end{proof}

\begin{example}{Independence in $\R^3$}{ex:independence-R3}
\begin{enumerate}[label=(\alph*)]
  \item The canonical basis vectors $\vece{1}, \vece{2}, \vece{3}$ are
    linearly independent: if $a\,\vece{1}+b\,\vece{2}+c\,\vece{3}=\vzero$,
    then $(a,b,c)=(0,0,0)$.
  \item The vectors $(1,1,0)$, $(1,0,1)$, $(0,1,1)$ are linearly
    independent. Indeed,
    $\alpha(1,1,0)+\beta(1,0,1)+\gamma(0,1,1)=\vzero$ gives
    \[
      \alpha+\beta=0,\quad \alpha+\gamma=0,\quad \beta+\gamma=0,
    \]
    which has the unique solution $\alpha=\beta=\gamma=0$.
  \item The vectors $(1,2,3)$, $(4,5,6)$, $(7,8,9)$ are linearly
    \emph{dependent}: $(1,2,3) - 2(4,5,6) + (7,8,9) = (0,0,0)$.
\end{enumerate}
\end{example}

\begin{example}{Independence in $\K_n[X]$}{ex:independence-polynomials}
The polynomials $1,X,X^2,\ldots,X^n$ are linearly independent in
$\K[X]$: if $\alpha_0 + \alpha_1 X + \cdots + \alpha_n X^n = 0$ (the
zero polynomial), then all $\alpha_i = 0$ by comparing coefficients.
\end{example}


% ============================================================
\section{Bases}\label{sec:bases}
% ============================================================

\begin{definition}{Basis}{def:basis}
\index{basis!definition}
A family $\cB = (v_1,\ldots,v_n)$ of vectors in $E$ is a \emph{basis}
of $E$ if it is both linearly independent and Spanning: $\cB$ is
linearly independent and $\Span(\cB) = E$.
\end{definition}

\begin{theorem}{Characterization of a basis}{thm:basis-char}
\index{basis!characterization}
Let $\cB = (v_1,\ldots,v_n)$ be a family in $E$. The following are
equivalent:
\begin{enumerate}[label=(\roman*)]
  \item $\cB$ is a basis of $E$.
  \item Every $v\in E$ can be written \emph{uniquely} as
    $v = \lambda_1 v_1 + \cdots + \lambda_n v_n$.
  \item $\cB$ is a maximal linearly independent family.
  \item $\cB$ is a minimal Spanning family.
\end{enumerate}
\end{theorem}

\begin{proof}
\textbf{(i)$\Rightarrow$(ii):} Since $\cB$ Spans $E$, at least one
representation exists. Suppose
$v = \sum_i\lambda_i v_i = \sum_i\mu_i v_i$. Then
$\sum_i(\lambda_i-\mu_i)v_i = \vzero$, and independence gives
$\lambda_i = \mu_i$ for all $i$.

\textbf{(ii)$\Rightarrow$(i):} Existence of a representation means $\cB$
Spans $E$. If $\sum_i\lambda_i v_i = \vzero$, then
$\sum_i\lambda_i v_i = \sum_i 0\cdot v_i$ are two representations of
$\vzero$; uniqueness gives $\lambda_i = 0$ for all $i$.

\textbf{(i)$\Rightarrow$(iii):} If $\cB$ is a basis and we add any
$w\in E$, then $w = \sum_i\lambda_i v_i$, so the enlarged family
contains $w - \sum_i\lambda_i v_i = \vzero$ as a nontrivial relation;
hence it is dependent. So $\cB$ is maximal independent.

\textbf{(iii)$\Rightarrow$(i):} If $\cB$ does not Span $E$, there exists
$w\notin\Span(\cB)$, and $(v_1,\ldots,v_n,w)$ is independent by

ef{prop:prop:li-properties}(iii), contradicting maximality.

\textbf{(i)$\Rightarrow$(iv):} If we remove $v_j$, we claim the
remaining family does not Span $E$. Indeed, $v_j$ cannot be written as a
linear combination of the other $v_i$ (by independence), so $v_j$ is not
in the Span of the reduced family.

\textbf{(iv)$\Rightarrow$(i):} $\cB$ Spans $E$ by assumption. If $\cB$
is dependent, some $v_j$ is a combination of the others (by

ef{prop:prop:ld-char}), so removing $v_j$ still gives a Spanning family,
contradicting minimality.
\end{proof}

\begin{definition}{Coordinates}{def:coordinates}
\index{coordinates}
If $\cB = (v_1,\ldots,v_n)$ is a basis of $E$ and $v = \sum_i\lambda_iv_i$,
the uniquely determined scalars $\lambda_1,\ldots,\lambda_n$ are the
\emph{coordinates} (or \emph{components}) of $v$ relative to $\cB$. The
column vector
\[
  [v]_\cB = \begin{pmatrix}\lambda_1\\\vdots\\\lambda_n\end{pmatrix}
  \in\K^n
\]
is called the \emph{coordinate vector} of $v$ in basis $\cB$.
\end{definition}

\begin{example}{Canonical basis of $\K^n$}{ex:canonical-basis}
\index{canonical basis!of $\K^n$}
The \emph{canonical} (or \emph{standard}) basis of $\K^n$ is
$\cE = (\vece{1},\ldots,\vece{n})$ where $\vece{i}$ is the $n$-tuple
with a $1$ in position $i$ and $0$'s elsewhere. The coordinates of
$(x_1,\ldots,x_n)$ in $\cE$ are simply $x_1,\ldots,x_n$.
\end{example}

\begin{example}{Canonical basis of $\K_n[X]$}{ex:canonical-basis-poly}
\index{canonical basis!of polynomial space}
The family $(1,X,X^2,\ldots,X^n)$ is a basis of $\K_n[X]$: every
polynomial of degree $\leq n$ is a unique linear combination of these
monomials. In this basis, the coordinates of a polynomial are its
coefficients.
\end{example}

\begin{example}{Canonical basis of $\Mat_{n,p}(\K)$}{ex:canonical-basis-matrix}
\index{canonical basis!of matrix space}
The family $(E_{ij})_{1\leq i\leq n,\,1\leq j\leq p}$, where $E_{ij}$
is the matrix with $1$ in position $(i,j)$ and $0$'s elsewhere, is a
basis of $\Mat_{n,p}(\K)$.
\end{example}


% ============================================================
\section{Dimension}\label{sec:dimension}
% ============================================================

The notion of dimension captures the ``size'' of a vector space. Its
definition rests on a fundamental theorem about the relationship between
independent families and Spanning families.

\begin{theorem}{Steinitz exchange lemma}{thm:steinitz}
\index{Steinitz exchange lemma}
Let $E$ be a $\K$-vector space. If $(v_1,\ldots,v_p)$ is linearly
independent and $(w_1,\ldots,w_q)$ Spans $E$, then $p\leq q$.
\end{theorem}

\begin{proof}
We proceed by induction on $p$.

\textbf{Base case ($p=0$):} The empty family is independent and $0\leq q$.

\textbf{Inductive step:} Assume the result holds for all independent
families of size $p-1$. Suppose $(v_1,\ldots,v_p)$ is independent and
$(w_1,\ldots,w_q)$ Spans $E$.

Since $(w_1,\ldots,w_q)$ Spans $E$, we can write
\[
  v_p = \alpha_1 w_1 + \cdots + \alpha_q w_q
\]
for some $\alpha_i\in\K$. Not all $\alpha_i$ can be zero (since
$v_p\neq\vzero$). By reindexing, assume $\alpha_q\neq 0$. Then
\[
  w_q = \alpha_q^{-1}\Bigl(v_p - \sum_{i=1}^{q-1}\alpha_i w_i\Bigr).
\]
Therefore every vector in $E$ that was a linear combination of
$w_1,\ldots,w_q$ can also be expressed using $w_1,\ldots,w_{q-1},v_p$.
That is, $(w_1,\ldots,w_{q-1},v_p)$ still Spans $E$.

Now $(v_1,\ldots,v_{p-1})$ is independent (as a subfamily of an
independent family), and $(w_1,\ldots,w_{q-1},v_p)$ Spans $E$. We claim
$(v_1,\ldots,v_{p-1})$ remains independent in $E$, so by the induction
hypothesis, $p-1\leq q-1+1 = q$... but we need to be more careful.

We refine the argument. Since $(w_1,\ldots,w_{q-1},v_p)$ Spans $E$,
each $v_i$ ($1\leq i\leq p-1$) can be written as a combination of
$w_1,\ldots,w_{q-1},v_p$. Since $(v_1,\ldots,v_p)$ is independent,
$v_1,\ldots,v_{p-1}$ are not in $\Span(v_p)$ and hence depend
nontrivially on $w_1,\ldots,w_{q-1}$. We may thus view
$(v_1,\ldots,v_{p-1})$ as a linearly independent family in the space
$\Span(w_1,\ldots,w_{q-1},v_p)$.

Apply the exchange process again: write $v_{p-1}$ as a combination of
$w_1,\ldots,w_{q-1},v_p$. The coefficient of some $w_j$ must be nonzero
(otherwise $v_{p-1}\in\Span(v_p)$, contradicting independence). Swap
$w_j$ for $v_{p-1}$.

Continuing this process $p$ times, we replace $p$ of the $w_i$'s by
$v_1,\ldots,v_p$. At each step, at least one $w_i$ must be available
for replacement (its coefficient must be nonzero). This is possible only
if $p\leq q$.
\end{proof}

\begin{corollary}{Invariance of basis size}{cor:basis-size}
\index{basis!invariance of size}
If $E$ has a finite basis, then every basis of $E$ has the same number of
elements.
\end{corollary}

\begin{proof}
Let $\cB = (v_1,\ldots,v_n)$ and $\cB' = (w_1,\ldots,w_m)$ be two
bases. Since $\cB$ is independent and $\cB'$ Spans $E$, the Steinitz
lemma gives $n\leq m$. Symmetrically, $m\leq n$. Hence $n=m$.
\end{proof}

\begin{definition}{Dimension}{def:dimension}
\index{dimension}
A $\K$-vector space $E$ is \emph{finite-dimensional} if it has a finite
Spanning set. In this case, the \emph{dimension} of $E$, denoted
$\dim_\K E$ (or simply $\dim E$), is the common cardinality of all its
bases.

By convention, $\dim\set{\vzero} = 0$.

If $E$ has no finite Spanning set, it is \emph{infinite-dimensional}.
\end{definition}

\begin{example}{Dimensions of standard spaces}{ex:dimensions}
\begin{enumerate}[label=(\alph*)]
  \item $\dim_\K \K^n = n$ (the canonical basis has $n$ elements).
  \item $\dim_\K \K_n[X] = n+1$ (the basis $1,X,\ldots,X^n$ has $n+1$
    elements).
  \item $\dim_\K \Mat_{n,p}(\K) = np$.
  \item $\K[X]$ is infinite-dimensional: the family
    $(1,X,X^2,\ldots)$ is independent and Spans $\K[X]$, but it is
    not finite.
\end{enumerate}
\end{example}

\begin{theorem}{Existence of a basis (finite-dimensional case)}{thm:basis-existence}
\index{basis!existence}
Every finite-dimensional vector space $E\neq\set{\vzero}$ has a basis.
More precisely:
\begin{enumerate}[label=(\roman*)]
  \item Every Spanning family contains a basis (one can extract a basis
    from it).
  \item Every linearly independent family can be extended to a basis.
\end{enumerate}
\end{theorem}

\begin{proof}
\begin{enumerate}[label=(\roman*)]
  \item Let $\cF = (w_1,\ldots,w_q)$ Span $E$. If $\cF$ is independent,
    it is a basis. If not, some $w_j$ is a combination of the others;
    remove it. The resulting family still Spans $E$ (any combination
    involving $w_j$ can be rewritten using the others). Repeat until the
    family is independent. The process terminates since we cannot reduce
    below a single vector (assuming $E\neq\set{\vzero}$).

  \item Let $(v_1,\ldots,v_p)$ be independent. If it Spans $E$, it is a
    basis. If not, there exists $w\notin\Span(v_1,\ldots,v_p)$, and
    $(v_1,\ldots,v_p,w)$ is independent by
    
ef{prop:prop:li-properties}(iii). Repeat. By the Steinitz lemma, the
    size of an independent family is bounded by any Spanning family's
    size, so the process terminates. \qedhere
\end{enumerate}
\end{proof}

\begin{theorem}{Existence of complements}{thm:complement-existence}
\index{complementary subspace!existence}
Let $E$ be a finite-dimensional $\K$-vector space and $F$ a subspace of
$E$. Then $F$ has at least one complement in $E$: there exists a
subspace $G$ such that $E = F\oplus G$.
\end{theorem}

\begin{proof}
Let $(v_1,\ldots,v_p)$ be a basis of $F$. This is a linearly independent
family in $E$; by 
ef{thm:thm:basis-existence}(ii), extend it to a basis
$(v_1,\ldots,v_p,w_1,\ldots,w_q)$ of $E$. Set
$G = \Span(w_1,\ldots,w_q)$.

\textbf{$E = F+G$:} Any $x\in E$ can be written as
$x = \sum_i\alpha_iv_i + \sum_j\beta_jw_j\in F+G$.

\textbf{$F\cap G = \set{\vzero}$:} If $x\in F\cap G$, write
$x = \sum_i\alpha_iv_i = \sum_j\beta_jw_j$. Then
$\sum_i\alpha_iv_i - \sum_j\beta_jw_j = \vzero$, and since
$(v_1,\ldots,v_p,w_1,\ldots,w_q)$ is a basis (hence independent), all
coefficients vanish: $x = \vzero$.
\end{proof}


% ============================================================
\section{Dimension formulas}\label{sec:dimension-formulas}
% ============================================================

\begin{proposition}{Dimension of a subspace}{prop:dim-subspace}
\index{dimension!of a subspace}
Let $E$ be a finite-dimensional $\K$-vector space and $F$ a subspace of
$E$. Then:
\begin{enumerate}[label=(\roman*)]
  \item $F$ is finite-dimensional and $\dim F\leq\dim E$.
  \item $\dim F = \dim E$ if and only if $F = E$.
\end{enumerate}
\end{proposition}

\begin{proof}
\begin{enumerate}[label=(\roman*)]
  \item Any independent family in $F$ is also independent in $E$, so its
    size is at most $\dim E$ (by the Steinitz lemma). Take a maximal
    independent family in $F$; it is a basis of $F$ (if it did not Span
    $F$, it could be enlarged, contradicting maximality). Hence $F$ is
    finite-dimensional and $\dim F\leq\dim E$.

  \item If $\dim F = \dim E = n$, let $(v_1,\ldots,v_n)$ be a basis
    of $F$. This is an independent family of size $n$ in $E$, hence a
    basis of $E$ (a maximal independent family, since no independent
    family in $E$ has more than $\dim E = n$ elements). Thus
    $F = \Span(v_1,\ldots,v_n) = E$.

    Conversely, $F = E$ trivially implies $\dim F = \dim E$. \qedhere
\end{enumerate}
\end{proof}

\begin{proposition}{Dimension of a direct sum}{prop:dim-direct-sum}
\index{dimension!of a direct sum}
If $E = F\oplus G$, then $\dim E = \dim F + \dim G$.
\end{proposition}

\begin{proof}
Let $(v_1,\ldots,v_p)$ be a basis of $F$ and $(w_1,\ldots,w_q)$ a basis
of $G$. We show $(v_1,\ldots,v_p,w_1,\ldots,w_q)$ is a basis of $E$.

\textbf{Spanning:} Any $x\in E$ can be written $x = u + v$ with $u\in F$,
$v\in G$. Expanding $u$ and $v$ in their respective bases gives $x$ as a
combination of $v_1,\ldots,v_p,w_1,\ldots,w_q$.

\textbf{Independence:} Suppose
$\sum_i\alpha_i v_i + \sum_j\beta_j w_j = \vzero$. Then
$\sum_i\alpha_iv_i = -\sum_j\beta_jw_j\in F\cap G = \set{\vzero}$.
Hence $\sum_i\alpha_iv_i = \vzero$ and $\sum_j\beta_jw_j = \vzero$,
which by independence of each basis gives all $\alpha_i = 0$ and all
$\beta_j = 0$.

Therefore the combined family is a basis of $E$ with $p+q$ elements.
\end{proof}

\begin{corollary}{Dimension of a complement}{cor:dim-complement}
If $F$ is a subspace of a finite-dimensional space $E$, then any
complement $G$ of $F$ satisfies $\dim G = \dim E - \dim F$.
\end{corollary}


% ============================================================
\section{The Grassmann formula}\label{sec:grassmann}
% ============================================================

\begin{theorem}{Grassmann formula (dimension of a sum)}{thm:grassmann}
\index{Grassmann formula}
Let $E$ be a finite-dimensional $\K$-vector space and let $F,G$ be
subspaces of $E$. Then
\[
  \dim(F+G) = \dim F + \dim G - \dim(F\cap G).
\]
\end{theorem}

\begin{proof}
Let $p = \dim F$, $q = \dim G$, and $r = \dim(F\cap G)$.

\textbf{Step 1.} Let $(u_1,\ldots,u_r)$ be a basis of $F\cap G$.

\textbf{Step 2.} Since $F\cap G\subseteq F$, we can extend
$(u_1,\ldots,u_r)$ to a basis of $F$:
\[
  (u_1,\ldots,u_r,\,v_1,\ldots,v_s) \quad\text{is a basis of } F,
\]
where $s = p - r$.

\textbf{Step 3.} Similarly, extend $(u_1,\ldots,u_r)$ to a basis of $G$:
\[
  (u_1,\ldots,u_r,\,w_1,\ldots,w_t) \quad\text{is a basis of } G,
\]
where $t = q - r$.

\textbf{Step 4.} We claim that
\[
  \cB = (u_1,\ldots,u_r,\,v_1,\ldots,v_s,\,w_1,\ldots,w_t)
\]
is a basis of $F+G$.

\textit{Spanning:} Any $x\in F+G$ is $x = f+g$ with $f\in F$ and
$g\in G$. Writing $f$ in the basis of $F$ and $g$ in the basis of $G$:
\[
  x = \Bigl(\sum_{i=1}^r \alpha_i u_i + \sum_{j=1}^s \beta_j v_j\Bigr)
    + \Bigl(\sum_{i=1}^r \gamma_i u_i + \sum_{k=1}^t \delta_k w_k\Bigr)
  = \sum_{i=1}^r(\alpha_i+\gamma_i)u_i
    + \sum_{j=1}^s\beta_j v_j
    + \sum_{k=1}^t\delta_k w_k.
\]
So $x\in\Span(\cB)$.

\textit{Independence:} Suppose
\begin{equation}\label{eq:grassmann-indep}
  \sum_{i=1}^r \alpha_i u_i + \sum_{j=1}^s \beta_j v_j
  + \sum_{k=1}^t \delta_k w_k = \vzero.
\end{equation}
Set $h \deff \sum_{k=1}^t\delta_k w_k =
-\sum_{i=1}^r\alpha_i u_i - \sum_{j=1}^s\beta_j v_j$.

The right-hand side lies in $F$ (since $u_i,v_j\in F$). The left-hand
side lies in $G$ (since $w_k\in G$). Hence $h\in F\cap G$.

Since $(u_1,\ldots,u_r)$ is a basis of $F\cap G$, we can write
$h = \sum_{i=1}^r\mu_i u_i$ for some $\mu_i\in\K$. But also
$h = \sum_{k=1}^t\delta_k w_k$, so
\[
  \sum_{k=1}^t\delta_k w_k - \sum_{i=1}^r\mu_i u_i = \vzero.
\]
Since $(u_1,\ldots,u_r,w_1,\ldots,w_t)$ is a basis of $G$ (hence
independent), all $\delta_k = 0$ and all $\mu_i = 0$. In particular,
$h = \vzero$.

Returning to~\eqref{eq:grassmann-indep} with all $\delta_k=0$:
$\sum_i\alpha_iu_i + \sum_j\beta_jv_j = \vzero$. Since
$(u_1,\ldots,u_r,v_1,\ldots,v_s)$ is a basis of $F$ (hence independent),
all $\alpha_i = 0$ and all $\beta_j = 0$.

\textbf{Conclusion.} $\cB$ is a basis of $F+G$ with
$r + s + t = r + (p-r) + (q-r) = p+q-r$ elements, so
\[
  \dim(F+G) = p + q - r = \dim F + \dim G - \dim(F\cap G). \qedhere
\]
\end{proof}

\begin{corollary}{Criterion for direct sum via dimension}{cor:direct-sum-dim}
\index{direct sum!dimension criterion}
Let $F$ and $G$ be subspaces of a finite-dimensional space $E$. Then
$E = F\oplus G$ if and only if $E = F+G$ and $\dim E = \dim F+\dim G$.
Equivalently, any two of the following three conditions imply the third:
\begin{enumerate}[label=(\roman*)]
  \item $E = F+G$,
  \item $F\cap G = \set{\vzero}$,
  \item $\dim E = \dim F + \dim G$.
\end{enumerate}
\end{corollary}

\begin{proof}
By the Grassmann formula, $\dim(F+G) = \dim F + \dim G - \dim(F\cap G)$.
Condition~(ii) says $\dim(F\cap G)=0$. Condition~(i) says
$\dim(F+G)=\dim E$. Condition~(iii) is
$\dim E = \dim F + \dim G$.

If (i) and (ii) hold: $\dim E = \dim(F+G) = \dim F+\dim G$, giving (iii).

If (i) and (iii) hold: $\dim F+\dim G = \dim E = \dim(F+G) =
\dim F+\dim G-\dim(F\cap G)$, so $\dim(F\cap G)=0$, giving (ii).

If (ii) and (iii) hold: $\dim(F+G) = \dim F+\dim G = \dim E$, and since
$F+G\subseteq E$ is a subspace with the same dimension as $E$,

ef{prop:prop:dim-subspace}(ii) gives $F+G=E$, i.e.\ (i).
\end{proof}


% ============================================================
\section{Rank of a family of vectors}\label{sec:rank}
% ============================================================

\begin{definition}{Rank of a family}{def:rank-family}
\index{rank!of a family of vectors}
Let $(v_1,\ldots,v_p)$ be a family of vectors in a $\K$-vector space
$E$. The \emph{rank} of this family is
\[
  \rank(v_1,\ldots,v_p) \deff \dim\Span(v_1,\ldots,v_p).
\]
\end{definition}

\begin{proposition}{Properties of rank}{prop:rank-properties}
Let $\cF = (v_1,\ldots,v_p)$ be a family of vectors in $E$.
\begin{enumerate}[label=(\roman*)]
  \item $0\leq\rank(\cF)\leq p$.
  \item $\rank(\cF) = p$ if and only if $\cF$ is linearly independent.
  \item If $E$ is finite-dimensional, $\cF$ Spans $E$ if and only if
    $\rank(\cF) = \dim E$.
  \item $\cF$ is a basis of $E$ if and only if
    $\rank(\cF) = p = \dim E$.
\end{enumerate}
\end{proposition}

\begin{proof}
\begin{enumerate}[label=(\roman*)]
  \item $\Span(\cF)$ is a subspace Spanned by $p$ vectors, so
    $\dim\Span(\cF)\leq p$. The lower bound is clear.

  \item $\rank(\cF) = p$ means $\Span(\cF)$ has a basis of size $p$.
    Since $\cF$ Spans $\Span(\cF)$, we can extract a basis of size
    $\dim\Span(\cF)=p$ from $\cF$; but $\cF$ already has $p$ elements,
    so it must be the basis itself---hence $\cF$ is independent.
    Conversely, if $\cF$ is independent, it is a basis of $\Span(\cF)$,
    so $\rank(\cF)=p$.

  \item $\cF$ Spans $E$ iff $\Span(\cF) = E$ iff
    $\dim\Span(\cF)=\dim E$ (by 
ef{prop:prop:dim-subspace}(ii)).

  \item Combines (ii) and (iii). \qedhere
\end{enumerate}
\end{proof}

\begin{remark}{Rank of the columns of a matrix}{rem:rank-matrix}
\index{rank!of a matrix}
If $A\in\Mat_{n,p}(\K)$ has columns $C_1,\ldots,C_p\in\K^n$, then
$\rank(C_1,\ldots,C_p) = \dim\Span(C_1,\ldots,C_p)$. This number is
called the \emph{column rank} of $A$. It equals the \emph{row rank} and
is simply called the \emph{rank} of $A$, denoted $\rank(A)$. The proof
of this equality will be given in 
ef{ch:linear-maps}.
\end{remark}


% ============================================================
\section{Applications}\label{sec:vs-applications}
% ============================================================

\subsection{Solution spaces of homogeneous linear
systems}\label{subsec:solution-spaces}

\begin{proposition}{Dimension of the solution space}{prop:solution-space-dim}
\index{solution space!dimension}
Let $A\in\Mat_{n,p}(\K)$ with $\rank(A)=r$. The solution space
$\cS = \setcond{X\in\K^p}{AX=\vzero}$ is a subspace of $\K^p$ of
dimension $p - r$.
\end{proposition}

\begin{proof}
That $\cS$ is a subspace was shown in 
ef{ex:ex:homogeneous-system}. The
dimension formula $\dim\cS = p - r$ is the \emph{rank--nullity theorem},
which will be proved rigorously in 
ef{ch:linear-maps} once the
necessary language of linear maps is available. For now, we note that
Gaussian elimination produces exactly $p - r$ free variables, each of
which gives rise to one basis vector of $\cS$.
\end{proof}

\begin{example}{A concrete system}{ex:concrete-system}
Consider the homogeneous system in $\R^4$:
\[
  \begin{cases}
    x_1 + 2x_2 - x_3 + x_4 = 0, \\
    2x_1 + 4x_2 + x_3 - 2x_4 = 0.
  \end{cases}
\]
The coefficient matrix has rank $2$ (since the two rows are not
proportional after reduction). Hence $\dim\cS = 4-2 = 2$.

Row reduction gives $x_1 = -2x_2 + x_3 - x_4$ and $x_3 = \frac{4}{3}x_4$
(after back-substitution). Setting the free variables $x_2$ and $x_4$ to
$(1,0)$ and $(0,1)$ respectively yields a basis of $\cS$.
\end{example}

\subsection{Polynomial spaces}\label{subsec:polynomial-applications}

\begin{example}{Even and odd polynomials}{ex:even-odd-poly}
\index{polynomial space!even and odd}
Let $E_n = \K_n[X]$ (polynomials of degree $\leq n$). Define
\begin{align*}
  F &= \setcond{P\in E_n}{P(-X)=P(X)} &&\text{(even polynomials)}, \\
  G &= \setcond{P\in E_n}{P(-X)=-P(X)} &&\text{(odd polynomials)}.
\end{align*}
Then $F$ and $G$ are subspaces of $E_n$ and $E_n = F\oplus G$.

Indeed, every polynomial $P$ can be uniquely decomposed as
\[
  P(X) = \underbrace{\frac{P(X)+P(-X)}{2}}_{\in F}
        + \underbrace{\frac{P(X)-P(-X)}{2}}_{\in G}.
\]
A basis of $F$ is $(1,X^2,X^4,\ldots)$ (even powers up to $n$) and a
basis of $G$ is $(X,X^3,X^5,\ldots)$ (odd powers up to $n$). The
dimensions satisfy
\[
  \dim F + \dim G = \left\lfloor\tfrac{n}{2}\right\rfloor+1
  + \left\lfloor\tfrac{n+1}{2}\right\rfloor = n+1 = \dim E_n. \qedhere
\]
\end{example}

\begin{example}{Symmetric and skew-symmetric matrices}{ex:sym-skew}
\index{symmetric matrix}\index{skew-symmetric matrix}
Let $E = \Mat_n(\K)$ with $\Char(\K)\neq 2$. Define
\begin{align*}
  S &= \setcond{A\in E}{\transpose{A}=A}, \\
  A &= \setcond{A\in E}{\transpose{A}=-A}.
\end{align*}
Then $E = S\oplus A$, since every matrix $M$ decomposes uniquely as
\[
  M = \underbrace{\frac{M+\transpose{M}}{2}}_{\in S}
    + \underbrace{\frac{M-\transpose{M}}{2}}_{\in A}.
\]
We have $\dim S = \frac{n(n+1)}{2}$ and $\dim A = \frac{n(n-1)}{2}$,
and indeed $\frac{n(n+1)}{2}+\frac{n(n-1)}{2}=n^2=\dim E$.
\end{example}


% ============================================================
\section{Exercises}\label{sec:ch1-exercises}
% ============================================================

\begin{exercise}{Verifying vector space axioms \basic}{exo:verify-axioms}
Verify that the set $\R^+ = (0,\infty)$ with the operations
\[
  x\oplus y \deff xy, \qquad \lambda\odot x \deff x^\lambda,
\]
is an $\R$-vector space. What is the zero vector? What is the additive
inverse of $x$?
\end{exercise}

\begin{exercise}{Subspace verification \basic}{exo:subspace-check}
Determine which of the following subsets of $\R^3$ are subspaces. Justify
your answers.
\begin{enumerate}[label=(\alph*)]
  \item $S_1 = \setcond{(x,y,z)\in\R^3}{x+2y-z=0}$.
  \item $S_2 = \setcond{(x,y,z)\in\R^3}{x^2+y^2+z^2\leq 1}$.
  \item $S_3 = \setcond{(x,y,z)\in\R^3}{x=y=z}$.
  \item $S_4 = \setcond{(x,y,z)\in\R^3}{xy=0}$.
  \item $S_5 = \setcond{(x,y,z)\in\R^3}{x+y+z=1}$.
\end{enumerate}
\end{exercise}

\begin{exercise}{Subspaces of function spaces \basic}{exo:function-subspace}
Let $E = \R^\R$ be the vector space of all functions $\R\to\R$. Prove
that the following subsets are subspaces:
\begin{enumerate}[label=(\alph*)]
  \item The set of even functions $\setcond{f\in E}{f(-x)=f(x)
    \text{ for all } x}$.
  \item The set of functions vanishing at $0$:
    $\setcond{f\in E}{f(0)=0}$.
  \item The set of periodic functions with period $T>0$:
    $\setcond{f\in E}{f(x+T)=f(x) \text{ for all } x}$.
\end{enumerate}
\end{exercise}

\begin{exercise}{Linear independence in $\R^4$ \basic}{exo:independence-R4}
Determine whether the following vectors in $\R^4$ are linearly
independent:
\[
  v_1 = (1,0,1,0),\quad v_2 = (0,1,0,1),\quad v_3 = (1,1,1,1),\quad
  v_4 = (1,0,0,1).
\]
If they are dependent, find an explicit nontrivial linear relation.
\end{exercise}

\begin{exercise}{Extending to a basis \intermediate}{exo:extend-basis}
In $\R^4$, the vectors $v_1=(1,1,0,0)$ and $v_2=(0,1,1,0)$ are linearly
independent. Extend $\set{v_1,v_2}$ to a basis of $\R^4$ by adding
suitable vectors from the canonical basis.
\end{exercise}

\begin{exercise}{Sum and intersection \intermediate}{exo:sum-intersection}
In $\R^4$, let
\begin{align*}
  F &= \Span\bigl((1,0,1,0),\,(0,1,0,1)\bigr), \\
  G &= \Span\bigl((1,1,0,0),\,(0,0,1,1)\bigr).
\end{align*}
\begin{enumerate}[label=(\alph*)]
  \item Find $\dim F$, $\dim G$, and $\dim(F+G)$.
  \item Determine $F\cap G$ and verify the Grassmann formula.
  \item Is $F+G$ a direct sum?
\end{enumerate}
\end{exercise}

\begin{exercise}{Direct sum decomposition \intermediate}{exo:direct-sum-decomp}
Let $E = \Mat_2(\R)$. Define
\[
  F = \setcond{A\in E}{\transpose{A}=A}, \qquad
  G = \setcond{A\in E}{\transpose{A}=-A}.
\]
\begin{enumerate}[label=(\alph*)]
  \item Show that $E = F\oplus G$.
  \item Find bases for $F$ and $G$, and verify that
    $\dim F+\dim G=\dim E$.
\end{enumerate}
\end{exercise}

\begin{exercise}{Polynomial subspace \intermediate}{exo:poly-subspace}
In $\R_3[X]$ (polynomials of degree $\leq 3$), let
\[
  F = \setcond{P\in\R_3[X]}{P(0)=0 \text{ and } P(1)=0}.
\]
\begin{enumerate}[label=(\alph*)]
  \item Show that $F$ is a subspace of $\R_3[X]$.
  \item Find a basis of $F$ and determine $\dim F$.
  \item Find a complement of $F$ in $\R_3[X]$.
\end{enumerate}
\end{exercise}

\begin{exercise}{The Grassmann formula in practice \intermediate}{exo:grassmann-practice}
Let $E=\R^5$ and
\begin{align*}
  F &= \setcond{(x_1,x_2,x_3,x_4,x_5)\in\R^5}
    {x_1-x_2=0,\; x_3-x_4+x_5=0}, \\
  G &= \setcond{(x_1,x_2,x_3,x_4,x_5)\in\R^5}
    {x_1=0,\; x_2-x_3=0,\; x_4=0}.
\end{align*}
\begin{enumerate}[label=(\alph*)]
  \item Find bases for $F$ and $G$.
  \item Compute $\dim F$, $\dim G$, $\dim(F\cap G)$, and $\dim(F+G)$.
  \item Verify the Grassmann formula.
\end{enumerate}
\end{exercise}

\begin{exercise}{Abstract dimension argument \intermediate}{exo:abstract-dim}
Let $E$ be a vector space of dimension $n$ and let $F,G$ be subspaces
with $\dim F + \dim G > n$. Prove that $F\cap G\neq\set{\vzero}$.
\end{exercise}

\begin{exercise}{Characterizing bases via cardinality \challenging}{exo:basis-card}
Let $E$ be a $\K$-vector space with $\dim E = n$. Let
$\cF = (v_1,\ldots,v_n)$ be a family of $n$ vectors. Prove that the
following are equivalent:
\begin{enumerate}[label=(\roman*)]
  \item $\cF$ is a basis of $E$.
  \item $\cF$ is linearly independent.
  \item $\cF$ Spans $E$.
\end{enumerate}
(Note: In general, one needs \emph{both} independence and Spanning to
guarantee a basis; the point of this exercise is that when the cardinality
equals $\dim E$, either condition alone suffices.)
\end{exercise}

\begin{exercise}{Direct sum of three subspaces \challenging}{exo:direct-sum-three}
In $\R^3$, let
\begin{align*}
  F_1 &= \Span\bigl((1,0,0)\bigr), \\
  F_2 &= \Span\bigl((0,1,0)\bigr), \\
  F_3 &= \Span\bigl((1,1,1)\bigr).
\end{align*}
\begin{enumerate}[label=(\alph*)]
  \item Show that $\R^3 = F_1\oplus F_2\oplus F_3$.
  \item Exhibit three subspaces $G_1,G_2,G_3$ of $\R^3$ such that
    $G_i\cap G_j = \set{\vzero}$ for all $i\neq j$ and
    $G_1+G_2+G_3=\R^3$, but $\R^3\neq G_1\oplus G_2\oplus G_3$.

    (\emph{Hint:} pairwise trivial intersection does not imply direct sum
    for three or more subspaces.)
\end{enumerate}
\end{exercise}

\begin{exercise}{Infinite-dimensional spaces \challenging}{exo:infinite-dim}
\begin{enumerate}[label=(\alph*)]
  \item Prove that $\K[X]$ is infinite-dimensional by showing that the
    family $(1,X,X^2,\ldots)$ is linearly independent.
  \item Show that the function space $\R^\R$ is infinite-dimensional.
    (\emph{Hint:} Consider the family $(\exp_\lambda)_{\lambda\in\R}$
    where $\exp_\lambda(x)=e^{\lambda x}$, or use the inclusion
    $\K[X]\hookrightarrow\R^\R$ via the evaluation map.)
  \item Prove that every vector space has a basis (assuming the Axiom of
    Choice, or equivalently, Zorn's Lemma).
\end{enumerate}
\end{exercise}

\begin{exercise}{Grassmann formula for three subspaces \challenging}{exo:grassmann-three}
Let $F_1,F_2,F_3$ be subspaces of a finite-dimensional space $E$. Prove
the inequality
\[
  \dim(F_1+F_2+F_3) \geq \dim F_1+\dim F_2+\dim F_3
  -\dim(F_1\cap F_2)-\dim(F_1\cap F_3)-\dim(F_2\cap F_3).
\]
Give an example where strict inequality holds.
\end{exercise}


% ============================================================
\section*{Chapter summary}\label{sec:ch1-summary}
\addcontentsline{toc}{section}{Chapter summary}
% ============================================================

\begin{itemize}
  \item A \textbf{vector space} over a field $\K$ is a set $E$ with
    addition and scalar multiplication satisfying eight axioms
    \ref{ax:assoc}--\ref{ax:dist-field}. The central examples are
    $\K^n$, $\Mat_{n,p}(\K)$, polynomial spaces $\K_n[X]$ and $\K[X]$,
    and function spaces.

  \item A \textbf{subspace} is a nonempty subset closed under addition
    and scalar multiplication. The intersection of subspaces is a
    subspace; their union is generally not, and is replaced by their
    \textbf{sum}.

  \item The sum $E = F\oplus G$ is \textbf{direct} when
    $F\cap G = \set{\vzero}$, equivalently when every vector has a
    unique $F$-plus-$G$ decomposition. Every subspace of a
    finite-dimensional space admits a complement.

  \item A \textbf{linear combination} of $v_1,\ldots,v_n$ is a vector
    $\sum\lambda_iv_i$. The \textbf{Span} of a family is the subspace of
    all such combinations. A family is \textbf{linearly independent} if
    the only combination yielding $\vzero$ is the trivial one.

  \item A \textbf{basis} is a family that is both independent and
    Spanning, or equivalently, a family giving a \emph{unique}
    representation for every vector. The \textbf{Steinitz exchange lemma}
    ensures that all bases have the same size; this common size is the
    \textbf{dimension}.

  \item Key dimension formulas:
    \begin{itemize}
      \item $\dim(F\oplus G) = \dim F + \dim G$.
      \item \textbf{Grassmann formula:}
        $\dim(F+G)=\dim F+\dim G-\dim(F\cap G)$.
      \item $F\subseteq E$ and $\dim F=\dim E$ imply $F = E$.
    \end{itemize}

  \item The \textbf{rank} of a family of vectors is the dimension of its
    Span. It equals the maximal number of independent vectors in the
    family.

  \item \textbf{Applications:} The solution space of a homogeneous system
    $AX=\vzero$ is a subspace of $\K^p$ of dimension $p-\rank(A)$.
    Polynomial spaces and matrix spaces admit natural direct sum
    decompositions (e.g.\ even/odd polynomials, symmetric/skew-symmetric
    matrices).
\end{itemize}
